% Priority-one namespace: target law, hidden forest, and target-class
% parameters.  Algorithmic and lower-witness notation must not reuse the
% rendered letters below for a different global meaning.

% Positions, vocabulary, target law, and forest.
\newcommand{\PositionSet}{[N]}
\newcommand{\Vocab}{\mathcal V}
\newcommand{\VocabSize}{V}
\newcommand{\TargetLaw}{P}
\newcommand{\InstanceTargetLaw}[1]{P^{#1}}
\newcommand{\TargetForest}{F}
\newcommand{\TargetEdges}{E}
\newcommand{\TargetRoots}{\operatorname{Roots}(\TargetForest)}
\newcommand{\TokenOrder}{\prec}
\newcommand{\ParentOf}[1]{\operatorname{pa}(#1)}

% Local probabilistic objects of the target law.
\newcommand{\PosMarginal}[1]{\pi_{#1}}
\newcommand{\EndpointKernel}[2]{K_{#1\to #2}}
\newcommand{\BoundaryRow}[4]{\mu_{#1\to #2}(\cdot\mid #3,#4)}
\newcommand{\TailSet}[2]{\mathcal T_{#1}(#2)}
\newcommand{\DirectedResponse}[3]{\Omega_{#1\to #2}(#3)}

% Public finite target-class parameters.
% The response exponent is deliberately undecorated: it is part of the target
% class, not a derived bank exponent or an algorithmic tuning parameter.
\newcommand{\ResponseConstant}{L}
\newcommand{\ResponseExponent}{\alpha}
\newcommand{\EdgeSignal}{\omega}
\newcommand{\FrequencyConstant}{C}

% Public asymptotic target-class exponents and hidden marginal ranks.
\newcommand{\VocabGrowthExponent}{\nu}
\newcommand{\RankTailExponent}{s}
\newcommand{\EdgeSignalExponent}{r}
\newcommand{\RankedToken}[2]{a_{#1,#2}}
